\chapter{Introduction}
\label{ch:introduction}

Where people look in a scene is not random. Human subjects converge on faces, text and salient objects
\citep{judd2009mit1003}. \emph{Where} attention lands is one question; in \emph{what order} it
lands is another. Computational models of gaze predict these from the image, fitted to eye-tracking
recordings of people looking at photographs. The earlier ones predicted only where, marking the regions
of an image that draw fixations. DeepGaze~III \citep{kummerer2022deepgaze3}, \dg{} hereafter, addresses
the second, the free-viewing \emph{scanpath} (\figref{intro-scanpath}). It predicts a probability density
over the next fixation given the image and the recent history, and is evaluated at the level of that
distribution rather than by matching individual sampled paths \citep{kummerer2021state}.

\begin{figure}[tbp]
  \centering
  \includegraphics[width=0.72\textwidth]{foveation_mit1003/intro_scanpath_stim0091.png}
  \caption[A human scanpath on an MIT1003 image]{One subject's fixations on an MIT1003 image \citep{judd2009mit1003}, numbered in viewing
    order (the first in white): the boarding stairs, the airline name, then the people on the tarmac.}
  \label{fig:intro-scanpath}
\end{figure}

Human vision is foveated, with acuity high only in a
small region around the point of gaze, falling steeply with eccentricity. That is why the eyes move, bringing
each region of interest onto the high-acuity fovea, with the choice of where to look next made from a
low-resolution peripheral preview \citep{najemnik2005}. \dg{} instead reads the whole image at uniform
resolution, so its predictions can be driven by fine detail a foveated viewer could not perceive.
Whether \dg{} needs that detail is the question this thesis puts to it.

This thesis narrows that gap at the input. A gaze-contingent, space-variant foveation after
\citet{geislerperry1998} re-blurs the model's input around its current fixation before every forward pass
--- sharp at gaze, blurred in the periphery --- so the model must predict each saccade
from the degraded peripheral preview a human would have had. Two otherwise-identical models, one trained on sharp
input and one on foveated input, are compared on MIT1003 \citep{judd2009mit1003} under image-stratified
cross-validation, alongside a third arm that pins the same blur to the image centre and so
separates gaze-contingency from the loss of peripheral resolution. At human acuity foveating the input costs
nothing the experiment can measure. A cost appears once the falloff is pushed past human
acuity, and grows as it steepens. Holding the same blur at the image centre costs more than letting it
follow the gaze, and the gap grows as the periphery is degraded further.
Alongside the foveation experiment, a per-image diagnostic shows
that \dg{}'s fit varies more from image to image than any arm differs from another. Future work sets out a
developmental question, whether a model's learning trajectory recapitulates the development of human
gaze across age.

\section{Applications}
\label{sec:intro-applications}

The mechanism and the measurement bear on problems outside the one asked here. \citet{geislerperry1998}
built the falloff this thesis borrows in order to send video over a low-bandwidth link, keeping detail
only where the eye was pointed. Gaze-tracked displays now do the same inside the renderer, drawing the
periphery of a rendered frame at reduced cost \citep{guenter2012foveated, patney2016foveated} and
steering what is left of the budget with a predicted saliency map \citep{gao2025foveatedpathtracing}. A
scanpath model gives the next fixation as well as the current one, so the region about to be looked at
can be drawn sharply before the eye reaches it. \chref{results} measures how the fit to human fixations
changes as peripheral detail is discarded, which is that budget question in another form.

The same problem arises with no human eye in the loop. A robot or a vehicle can carry a wide
low-resolution camera for the field of view and a narrow high-resolution one that is steered, which
\citet{wloka2018starfc} identify as the engineering counterpart of the retina's own arrangement.
Something has to decide where the narrow camera points, and a model that predicts the next fixation is a
candidate for that decision. \citet{zhang2020gazereview} survey the use of recorded human gaze in
training agents that face this choice.

Predicting where people look is already used in design. \citet{judd2009mit1003} name advertising among
the applications of fixation prediction, and \citet{wloka2018starfc} note that commercial tools predict
the first several fixations, the short window of attention an advertisement receives. A scanpath model
states the order in which the parts of a layout are reached, so a designer can ask whether the brand is
fixated early and not only whether it stands out. Interface layouts have been studied the same way
\citep{jiang2024eyeformer}.

\section{Structure of this thesis}

\chref{background} reviews foveated vision and the lineage of saliency and scanpath models leading to
\dg{}. \chref{methods} describes the host architecture and where the foveation attaches, then defines the
evaluation framework, the multi-subject and per-image diagnostics, and the gaze-contingent foveation
mechanism with its training protocol. \chref{results} reports the foveation experiment across the sharp
control, the gaze-contingent arms and the fixed-centre arms, then the per-image
analysis. \chref{discussion} interprets the finding and the per-image diagnostic, states the
limitations (among them what the dataset does not record about its subjects) and sets out the work the
mechanism enables. \chref{conclusion} states what the experiment settles and what it leaves open.
